% !PRINTED VIII
\weberpage{VIII}
sich im wesentlichen an die Abhandlung von \gesperrt{Dedekind} und mir im 92. Bande von Crelles Journal anschließt, aber durch Anwendung der Theorie der Funktionale, auf die ich im zweiten Bande die Theorie der algebraischen Zahlen gegründet habe, wie mir scheint, eine Vereinfachung erreicht.

\bigskip
\gesperrt{Straßburg}, im Mai 1908.

\begin{flushright}
\textbf{H. \gesperrt{Weber}.}
\end{flushright}
